\documentclass[12pt]{article}

\usepackage[utf8]{inputenc}
\usepackage[margin=1in]{geometry}
\usepackage{setspace}
\usepackage{url}

\doublespacing
\setlength{\emergencystretch}{3em}

\title{Housing Affordability and Belonging in Canada:\\
A Canadian News Archive Assignment}

\author{%
  Ma Toan Bach\\
  Student Number: 112708227\\
  LSO100KCC\\
  Canadian News Archive Assignment
}

\date{}

\begin{document}

\maketitle

\section*{Introduction}

Housing affordability has become one of the most important social issues in Canada today. The article I chose is ``Housing affordability in Canada is getting worse, federal records say,'' written by Craig Lord for \textit{The Canadian Press} and published by \textit{Global News} on September 12, 2025. The article discusses internal federal government documents showing that housing costs are making it harder for Canadians to find affordable places to live. It explains that vulnerable populations, lower-income households, middle-class Canadians, and newcomers are all affected by the lack of suitable and affordable housing (Lord, 2025). This issue matters because housing is connected to belonging, family stability, community, and social justice.

\section*{Summary}

The main idea of the article is that Canada's housing problem is becoming more serious, even though governments are trying to respond. According to Lord (2025), briefing materials prepared for Housing Minister Gregor Robertson described the housing situation as severe. The documents said that high housing costs are hurting the economy and making life more difficult for people who need affordable homes. Lower-income households and vulnerable people are struggling to meet basic housing needs, while middle-class Canadians are also having trouble buying homes and are staying in rental housing longer. This adds pressure to the rental market and can push rents even higher.

The article also explains that population growth, construction costs, and a shortage of non-market affordable housing have made the problem worse. Canada's stock of non-market affordable homes is lower than the OECD average, so there are not enough homes offered below regular market prices. The article also notes that average nightly homeless shelter use rose between 2020 and 2023, showing that the housing crisis is connected to homelessness (Lord, 2025). In response, the federal government planned to create Build Canada Homes to increase affordable homebuilding.

\section*{Connecting the News to Our Course}

This news story connects strongly to the course theme of belonging. In Canadian literature, belonging is often shown as the feeling of having a place where one is safe, accepted, and connected to others. Housing is a real-life part of that theme. A person may live in Canada, work in Canada, and contribute to society, but without a stable home, they may still feel insecure or excluded. A home gives people identity, routine, rest, and memory. Without stable housing, belonging becomes harder to experience.

This connects to Alice Munro's ``Boys and Girls'' because that story also shows how identity is shaped by the spaces people occupy and the roles they are allowed to have within them (Munro, 1968). In both the story and the housing article, ``home'' is not only a physical place. It is also connected to recognition, safety, and belonging.

The article also connects to family and community. When housing becomes too expensive, families may have to move far away from workplaces, schools, relatives, or cultural communities. Young adults may delay moving out, couples may delay having children, and newcomers may struggle to settle. These problems affect the whole community because people need stable housing to participate fully in society.

The Canada Mortgage and Housing Corporation report supports this issue by showing that even when new housing is being built, the type of housing matters. CMHC reported that housing starts increased in 2025, especially rental apartments and ``missing middle'' housing, but it also warned that homeownership supply is under pressure because condominium presales have weakened and some projects are being delayed or cancelled (Canada Mortgage and Housing Corporation, 2026). This shows why solving the crisis is complicated. Building more units is important, but those units must match what people need and can afford.

This issue also connects to social justice. Housing is a basic human need, but the people most affected are often those with fewer resources, including lower-income households, newcomers, and vulnerable Canadians. When housing is treated mainly as a market product, people with less money have fewer choices. This can deepen inequality through unstable housing, long commutes, overcrowding, or homelessness.

\section*{Reflection}

What surprised me most was how many different groups are affected by the housing crisis. I already knew that buying a home was difficult, especially in large Canadian cities, but the article showed that the issue goes beyond homeownership. Renters, newcomers, low-income families, and middle-class Canadians are all under pressure. I was also surprised by the connection between affordability and homelessness. The rise in shelter use shows how serious the crisis becomes when people do not have enough support.

This article strengthened my perspective because it helped me see housing as more than an economic issue. It is connected to dignity, identity, family life, and participation in society. When people cannot afford housing, they may feel like they do not fully belong in their own city or country. This connects to the course because many Canadian stories explore characters who search for home, safety, acceptance, and community. The housing crisis shows that these literary themes are part of real Canadian life.

\section*{Conclusion}

In conclusion, examining this news article helped me better understand how housing affordability affects Canadian society. The article shows that Canada's housing crisis is connected to economic pressure, population growth, construction challenges, and a lack of affordable non-market housing. More importantly, it shows that housing affects belonging, family stability, and community connection. By connecting the article to course themes, I learned that literature can help us think more deeply about real social issues.

\clearpage

\section*{References}

\hangindent=0.5in\noindent Canada Mortgage and Housing Corporation. (2026, March 11). \textit{Ownership supply under pressure as rentals dominate new construction}. \url{https://www.cmhc-schl.gc.ca/media-newsroom/news-releases/2026/spring-2026-housing-supply-report}

\vspace{1em}

\hangindent=0.5in\noindent Lord, C. (2025, September 12). Housing affordability in Canada is getting worse, federal records say. \textit{Global News}. \url{https://globalnews.ca/news/11417079/housing-affordability-canada-government-documents/}

\vspace{1em}

\hangindent=0.5in\noindent Munro, A. (1968). \textit{Boys and girls} [PDF]. \url{https://jerrywbrown.com/wp-content/uploads/2014/02/Boys-and-Girls-by-Alice-Munro.pdf}

\end{document}
